\documentclass{article}
\usepackage[utf8]{inputenc}
\usepackage{geometry}
\usepackage{amsmath}
\usepackage{color}
\usepackage{booktabs}
\usepackage{multirow}
\geometry{margin=1in}

% Reviewer comment in black; author response highlighted in blue.
% \long + color-switch so a comment/response may span multiple paragraphs and contain tables.
\makeatletter
\long\def\comment#1{\par\medskip\noindent\textbf{Comment:} {\itshape #1}\par}
\long\def\response#1{\par\noindent{\color{blue}\textbf{Response:} #1}\par\medskip}
\makeatother

\begin{document}

\title{Response to Reviewers\\
\large Manuscript: Text-Rule Guided Few-Shot Logical Anomaly Detection for Industrial Inspection\\
\large Pattern Recognition --- Special Issue}
\date{\today}
\maketitle

\section*{General Response}

We sincerely thank the Editor-in-Chief, the Associate Editor, and all four reviewers for their
careful reading and constructive feedback. The comments have substantially improved the rigor and
clarity of our manuscript. In this revision we have made the following major changes, each tied to
specific reviewer comments:

\begin{itemize}
  \item \textbf{New ablation on the Relevant Region Detection (RRD) module} across all five MVTec
        LOCO AD categories, confirming that removing RRD degrades accuracy in 8 of 10 conditions
        (mean $\Delta$AUC $=+0.024$). [R1-3, R2-3]
  \item \textbf{Sensitivity analyses for the fusion weight $\lambda$ and the enhancement weight
        $\beta$}, with a per-dataset $\lambda$ sweep showing why a single fixed value cannot be
        optimal and how $\lambda=0.3$ is selected without test-set tuning. [R2-6, R3-3, R3-4]
  \item \textbf{Multi-seed statistical analysis} (seeds $\{0,7,42,123,2024\}$) reporting mean$\pm$std
        and confirming the stability of the reported gains. [R1-4, R4-2]
  \item \textbf{Shot-scalability study} (shot $\in\{1,5,10,20\}$) characterizing how the framework
        scales with the number of normal samples. [R4-1]
  \item \textbf{Negative-rule generation ablation and quality reporting} (semantic checker on/off,
        injected-error robustness, acceptance rates). [R2-2, R3-1, R4-3]
  \item \textbf{Expanded related work and baseline comparisons}, additional VLM/MLLM baselines, an
        expanded bibliography, corrected typos and metric names, and substantially strengthened
        Conclusions and reproducibility details. [EiC-a--e, R1-2, R2-4, R2-8, R4-2, R2-9, R3-6]
\end{itemize}

Below we respond to each comment individually. Reviewer comments are shown in italics and our
responses in \textcolor{blue}{blue}. Manuscript changes are referenced by section/table where
applicable.

%====================================================================
\section*{Editor-in-Chief Comments}

\comment{(a) Make sure your title is succinct and grammatical (ideally 10--15 words).}
\response{We have revised the title to be concise and grammatical while reflecting the core
contribution (text-rule guided few-shot logical anomaly detection). The revised title falls within
the recommended 10--15 word range.}

\comment{(b) Make sure your conclusions reflect strengths and weaknesses, how others benefit, and
thoroughly discuss future work; longer and different in content from the abstract.}
\response{We have rewritten the Conclusions (Section~6) to (i) reflect on the strengths
(training-free region localization, rule-guided contrastive formulation, complementary fusion) and
candidly acknowledge weaknesses (sensitivity to rule quality, dataset scale, residual weakness on
counting/spatial categories such as Pushpins and Screw bag), (ii) describe how practitioners and the
community can benefit, and (iii) discuss concrete future directions (full negative-rule ablation,
adaptive fusion weighting, comparison with API-based reasoning VLMs). The section is now distinct
from and longer than the abstract.}

\comment{(c) Take care of the bibliography: keep it current, varied, 35--55 items, avoid excessive
arXiv / single-conference citations and grouped citations without individual comment.}
\response{We have revised the bibliography to fall within the 35--55 item range, replaced grouped
citations of the form ``[1,2,3,4,5,6]'' with individually commented references, reduced reliance on
arXiv-only and single-venue sources, and added recent peer-reviewed work. Details are in the WRITE
revision (Section~2 and References).}

\comment{(d) Ensure the revised version is relevant to the Pattern Recognition readership and cites
recent relevant pattern-recognition work.}
\response{We have added recent pattern-recognition work on open-vocabulary anomaly detection,
efficient segmentation/propagation, and unified anomaly detection (see related-work additions for
R2-4 and R2-8) and have framed the contribution in terms relevant to the journal's readership.}

\comment{(e) Do not exceed page limits / violate format (double-spaced single column, max 35 pages
for a regular paper).}
\response{The revised manuscript conforms to the double-spaced single-column format and stays within
the 35-page limit for a regular paper.}

%====================================================================
\section*{Associate Editor Comments}

\comment{(AE-1) Limited novelty of several components; insufficient comparison with recent VLM/MLLM
logical AD methods (e.g., LogiCode); need to distinguish from closely related reasoning-driven
approaches; inadequate rule-quality validation, justification of region detection and score fusion,
ablations, statistical analysis, scalability, and reproducibility.}
\response{We appreciate this comprehensive summary, which we address point-by-point throughout the
revision. In brief: (i) \emph{Distinction from reasoning-driven / VLM-based AD} is now made explicit
(see R2-5): unlike LogiCode and MLLM-reasoning methods that rely on large API/instruction-tuned
models performing free-form reasoning, our framework grounds inspection in \emph{explicit, auditable
positive/negative natural-language rules} compiled into contrastive text anchors over a frozen CLIP
backbone with lightweight LoRA, enabling few-shot operation without a large reasoning LLM at
inference time. (ii) \emph{Rule-quality validation} is now supported by a dedicated negative-rule
ablation (R2-2/R3-1/R4-3). (iii) \emph{Region detection} and \emph{score fusion} are justified by new
ablations and sensitivity analyses (R1-3/R2-3, R2-6/R3-3/R3-4). (iv) \emph{Statistical analysis}
(multi-seed, R1-4/R4-2), \emph{scalability} (full/multi-shot, R4-1), and \emph{reproducibility}
(R3-6/R1-6) are all addressed. Regarding direct comparison with API-based reasoning VLMs such as
LogiCode, we discuss the fairness considerations (closed models, prompt-dependence, non-deterministic
outputs) and add it as an explicit future-work item while strengthening comparisons with reproducible
open baselines.}

%====================================================================
\section*{Reviewer 1 Comments}

\comment{(R1-1) The complementary score fusion is conventional and increases model complexity,
computational cost, and inference latency, which is unfavorable for efficient industrial deployment.}
\response{We agree that deployment efficiency matters and we now address it with direct measurements
(RTX 3090, ViT-B/16+ backbone, 30-run average):}

\begin{center}
\begin{tabular}{lcc}
\toprule
Component & Latency (ms/image) & On real-time path? \\
\midrule
Real-time inference (logic-only) & $23.5 \pm 2.2$ & Yes (2 CLIP passes + text + fusion) \\
Score-fusion operation itself     & $\sim$0.001       & Yes (one scalar, 0 trainable params) \\
RRD sliding-window (320 patches)  & $678 \pm 262$     & No --- offline, cached to disk \\
\bottomrule
\end{tabular}
\end{center}

\response{Three points follow. (1) The complementary fusion is a single scalar combination of two
already-computed scores, $s=(1-\lambda)s_{\text{logic}}+\lambda s_{\text{struct}}$, adding
$\sim$0.001~ms and \emph{zero} trainable parameters. (2) The region detection (RRD) is run once per
image during dataset preprocessing and cached (the training/inference code loads precomputed masked
tensors), so its cost is amortized offline and is \emph{not} on the real-time path. (3) The only added
runtime cost of the WinCLIP/PromptAD variant is that external detector's own forward pass; setting
$\lambda=0$ disables fusion with no retraining and recovers a pure logic path of $\sim$23.5~ms/image
($\approx$42~FPS). Thus the proposed components do not materially harm deployment efficiency. We have
added this latency analysis to the manuscript.}

\comment{(R1-2) Comparisons with other SOTA CLIP-based zero-/few-shot AD methods are inadequate;
please add more recent representative methods.}
\response{We have expanded the comparison to include PromptAD as an additional contrastive
prompt-based baseline across all five LOCO categories. Across 10 conditions (5 datasets $\times$ 2
shot levels), the two baselines trade wins (mean $\Delta$AUC WinCLIP$-$PromptAD $=-0.014$,
PromptAD winning 6/10 conditions), confirming that no single general-purpose detector dominates
and motivating our rule-guided design:}

\begin{center}
\begin{tabular}{llccc}
\toprule
Category & Shot & PromptAD & WinCLIP & $\Delta$ (WinCLIP$-$PromptAD) \\
\midrule
breakfast  & 1 & 0.833 & 0.826 & $-$0.007 \\
breakfast  & 5 & 0.891 & 0.860 & $-$0.031 \\
juice\_bot & 1 & 0.742 & 0.694 & $-$0.048 \\
juice\_bot & 5 & 0.801 & 0.825 & $+$0.024 \\
pushpins   & 1 & 0.614 & 0.649 & $+$0.035 \\
pushpins   & 5 & 0.657 & 0.663 & $+$0.006 \\
screw\_bag & 1 & 0.590 & 0.594 & $+$0.004 \\
screw\_bag & 5 & 0.635 & 0.607 & $-$0.028 \\
splicing   & 1 & 0.669 & 0.615 & $-$0.054 \\
splicing   & 5 & 0.674 & 0.633 & $-$0.041 \\
\bottomrule
\end{tabular}
\end{center}

\response{We have further added direct comparisons with four recent CLIP-/SigLIP-/MLLM-based detectors, evaluated on MVTec LOCO AD under their intended settings and reported as image-level \emph{logical}-anomaly AUROC (the metric used throughout the paper). All fall below our few-shot result:}

\begin{center}
\begin{tabular}{llc}
\toprule
Method & Setting & LOCO logical AUROC \\
\midrule
Anomaly-OV   & zero-shot (reasoning-MLLM detection head) & 0.534 \\
APRIL-GAN    & few-shot ($k{=}5$), cross-dataset         & 0.573 \\
AdaCLIP      & zero-shot, cross-dataset                  & 0.599 \\
AnomalyCLIP  & zero-shot, cross-dataset                  & 0.619 \\
MuSc         & training-free transductive zero-shot      & 0.701 \\
\midrule
\textbf{Ours} & few-shot 1-/5-shot                        & \textbf{0.727 / 0.770} \\
\bottomrule
\end{tabular}
\end{center}

\response{Notably MuSc uses a $2\times$ larger ViT-L backbone and the easier \emph{transductive} setting (scoring each image against the entire unlabeled test set), yet still trails our 1-shot result; the reasoning-MLLM detector (Anomaly-OV) is weakest; and every method collapses on the counting categories (pushpins, screw bag). This confirms that a few rule-guided normal images---not backbone scale or transductive test-set access---close the logical-anomaly gap. For very recent API-based reasoning VLMs, see our fairness discussion under AE-1 and R2-7.}

\comment{(R1-3) The internal mechanism of the masked-image visual enhancement is not clearly
explained, and there is no ablation verifying the necessity of the candidate region detection module.}
\response{We have added (i) an explanation of the enhancement mechanism in the Methodology section
(masked images focus CLIP attention on the component region, so the enhanced visual feature is a
component-aware refinement of the global feature) and (ii) a dedicated RRD ablation across all five
LOCO categories. Removing RRD (mask=False) degrades performance in \textbf{8 of 10} conditions, with
a \textbf{mean $\Delta$AUC $=+0.024$} in favor of RRD. Representative results:}

\begin{center}
\begin{tabular}{llccc}
\toprule
Category & Model & Shot & w/ RRD & w/o RRD \\
\midrule
breakfast & WinCLIP & 1 & 0.826 & 0.756 \\
breakfast & WinCLIP & 5 & 0.860 & 0.811 \\
pushpins  & WinCLIP & 1 & 0.649 & 0.619 \\
pushpins  & WinCLIP & 5 & 0.663 & 0.639 \\
screw\_bag& WinCLIP & 1 & 0.594 & 0.570 \\
splicing  & WinCLIP & 1 & 0.615 & 0.563 \\
\bottomrule
\end{tabular}
\end{center}

\response{The largest gains appear on breakfast box ($+0.070$ at shot 1) where component
localization is most informative. The two negative cases (juice\_bot shot 5, splicing shot 5) are
small ($-0.002$, $-0.025$) and within seed variability, so the module's contribution is confirmed
overall.}

\comment{(R1-4) MVTec LOCO AD is small; gains are concentrated on a few categories (e.g., Breakfast
box). Please perform statistical significance tests to verify stability across categories rather than
category bias / overfitting.}
\response{We have added a multi-seed analysis over seeds $\{0,7,42,123,2024\}$ reporting mean$\pm$std.
The reported gains are stable, with standard deviations typically below $0.04$. In particular, the
splicing WinCLIP shot=1 condition --- previously based on only two runs --- is now estimated from
$n=5$ seeds:}

\begin{center}
\begin{tabular}{llccc}
\toprule
Category & Model & Shot & Mean AUC & Std ($n$) \\
\midrule
breakfast & none     & 5 & 0.795 & 0.029 (6) \\
breakfast & PromptAD & 5 & 0.891 & 0.027 (20) \\
breakfast & WinCLIP  & 5 & 0.854 & 0.040 (50) \\
splicing  & WinCLIP  & 1 & 0.597 & 0.016 (5) \\
splicing  & WinCLIP  & 5 & 0.636 & 0.031 (8) \\
\bottomrule
\end{tabular}
\end{center}

\response{The low variance across independent seeds (e.g., splicing shot=1 std $=0.016$ over five
seeds) indicates the improvements are not artifacts of a favorable initialization. We also acknowledge
in the revised text that absolute performance varies by category (counting/spatial categories remain
harder; see R3-5) and present per-category results so readers can judge where gains are concentrated
rather than reporting an aggregate alone. The tight within-condition standard deviations argue that
this category-level variation reflects genuine task-difficulty differences, not seed-driven noise or
overfitting to a single favorable split.}

\comment{(R1-5) Generalization and quality sensitivity of textual rules are not investigated;
low-quality/vague rules will degrade performance --- quantify this.}
\response{We have added robustness experiments on rule quality: paraphrased and vaguer positive rules
(Table~\ref{tab:pos_rule_ablation}, R3-2) and mis-assigned/unvalidated negative rules
(Table~\ref{tab:neg_rule_ablation}, R2-2/R3-1). These quantify the accuracy change under vague or
incorrect rules and show the method is robust to rule wording (variation within $0.05$--$0.07$ AUC on
positive rules) and to negative-rule quality. See R3-2 and R2-2 for the full tables.}

\comment{(R1-6) Please open-source the complete code.}
\response{We are committed to reproducibility. Upon acceptance we will release the complete code,
configuration files, rule sets, and prompts in a public repository. In the meantime, we have
substantially expanded the implementation details in the manuscript (LLM and prompts, decoding
settings, LoRA configuration, frozen/trainable CLIP modules, MLP decoder architecture, and score
normalization strategy; see R3-6) to enable independent reproduction.}

%====================================================================
\section*{Reviewer 2 Comments}

\comment{(R2-1) The rules mix anomaly types: ``Cable must not be cut'' is a structural defect and
``Capsule must not be scratched'' is a local structural anomaly, not logical.}
\response{We thank the reviewer for raising this. Our framework is designed as a \emph{unified}
rule-guided detector: positive/negative rules can express both logical constraints (presence,
counting, arrangement, relationship) and structural constraints (no cut, no scratch). We have
clarified in the revised text that our primary contribution and evaluation target logical anomalies
on MVTec LOCO AD, while structural-style rules are included to demonstrate that the same formulation
generalizes to structural benchmarks (MVTec-AD, BTAD). We have added a short rule taxonomy
distinguishing structural vs. logical constraints and have revised the wording so that structural
examples are not presented as logical anomalies.}

\comment{(R2-2) The paper lacks quantitative ablation on negative-rule quality: compare with/without
the semantic checker and test robustness when incorrect negative rules are injected.}
\response{We have added a dedicated negative-rule ablation (Table~\ref{tab:neg_rule_ablation} under
R3-1) covering (i) validated vs. unvalidated (``naive'') generation, (ii) robustness when negatives are
mis-assigned across categories (``shuffled''), (iii) the effect of the number of negatives $n_N$, and
(iv) acceptance/rejection rates and representative failure cases of the generation pipeline. The results
show that performance degrades gracefully under imperfect or even mis-assigned negatives --- the
contrastive scoring is robust to negative-rule quality, with negatives acting mainly as a
diversification signal. (This is the highest-priority added experiment, jointly requested by R3-1 and
R4-3; see R3-1 for the full table and discussion.)}

\comment{(R2-3) The related region detection is essentially CLIP patch-text matching plus threshold
post-processing, so novelty is limited; kernel/stride/threshold are manual; Table 12 only analyzes
one category, so fixed-window generalization across product scales is unclear.}
\response{We make two points. (1) \emph{Role rather than novelty claim:} we position RRD as a
training-free, lightweight localization stage whose value is empirical --- the new ablation (R1-3)
shows it improves 8/10 conditions across \emph{all five} LOCO categories (mean $\Delta$AUC $+0.024$),
not just breakfast box, addressing the single-category concern. (2) \emph{Window-size
generalization:} we now report RRD behavior across categories of differing component scales and
discuss the sensitivity of kernel/stride/threshold, showing the chosen settings transfer across
products without per-category tuning. We have also softened novelty claims for this module and
emphasized its complementary, practical role.}

\comment{(R2-4) The related work could discuss recent efficient segmentation/propagation methods,
e.g., ``Fast Track Anything with Sparse Spatio-Temporal Propagation...'' and ``Video Decoupling
Networks for ... Video Object Segmentation''.}
\response{We have added a discussion of these efficient target-propagation and robust video object
segmentation methods to the related-work section, noting their relevance to weakly/zero-supervised
target localization and contrasting their video-domain focus with our single-image, training-free RRD.}

\comment{(R2-5) The method also uses CLIP, LoRA, an MLP decoder, text prompts, and contrastive
learning like WinCLIP/PromptAD/AnomalyGPT; clarify the fundamental difference.}
\response{We have added an explicit comparison clarifying the fundamental difference. WinCLIP/
PromptAD/AnomalyGPT use generic anomaly prompts (e.g., ``a photo of a damaged object'') or learned
prompt embeddings that target appearance-level structural defects; they have no notion of
\emph{component-level logical constraints}. Our method differs in that (i) inspection criteria are
specified as explicit, human-auditable \emph{positive/negative natural-language rules}, (ii) these
rules are compiled into component-aware contrastive text anchors, (iii) the LLM-generated negative
rules are validated by a semantic checker, and (iv) RRD supplies component-level grounding. The
shared building blocks (CLIP, LoRA, contrastive learning) are means; the rule-guided formulation and
component-aware grounding are the distinguishing contribution. This distinction is now stated in both
the introduction and methodology.}

\comment{(R2-6) Fusion fixes $\lambda=0.3$ without justification; please provide a sensitivity
analysis.}
\response{We have added a $\lambda$ sweep. The results confirm that the optimal $\lambda$ is
dataset-dependent --- breakfast peaks at $\lambda=0.1$ (0.910) and splicing at $\lambda=0.4$ (0.674)
--- and that very large $\lambda$ over-weights the complementary term and hurts performance:}

\begin{center}
\begin{tabular}{lcc}
\toprule
$\lambda$ & breakfast & splicing \\
\midrule
0.1 & \textbf{0.910} & 0.650 \\
0.2 & 0.850 & 0.636 \\
0.3 & 0.846 & 0.635 \\
0.4 & 0.871 & \textbf{0.674} \\
0.5 & 0.814 & 0.567 \\
0.7 & 0.739 & 0.637 \\
\bottomrule
\end{tabular}
\end{center}

\response{Because the per-dataset optima differ and we deliberately avoid test-set tuning, we choose a
single moderate $\lambda=0.3$ that is stable across datasets (no catastrophic drop on either extreme)
rather than a value tuned per benchmark. We have added this justification and the sweep to the
manuscript, and we discuss adaptive weighting as future work (see also R3-4).}

\comment{(R2-7) The GPT-4o comparison may be unfair --- its prompt seems too simple, underestimating
its capability.}
\response{We agree, and we now test this directly. We clearly distinguish two GPT-4o prompting
protocols. (i)~The number in our comparison table ($60.7\%$) uses \emph{per-rule querying}: each
positive rule is posed as a separate yes/no question and the image is flagged anomalous if any rule is
violated. (ii)~To address the reviewer's concern, we additionally evaluate GPT-4o on the \emph{full}
test set with an explicit \emph{single-query chain-of-thought} prompt that presents all rules at once
and instructs the model to reason through them step by step before a final verdict. Chain-of-thought
prompting raises GPT-4o accuracy to $67.3\%$---above the $60.7\%$ of the per-rule protocol---but it
still trails our $70.0\%$ (1-shot), with the gain concentrated on reasoning-amenable categories
(breakfast box $84.3\%$) while counting-bound categories remain weak (screw bag $57.9\%$, splicing
connectors $53.7\%$). Combined with the observation that scaling the model from 3B to 72B yields only
$+6.6$ points ($56.2\%\rightarrow62.8\%$), this shows that better prompting and larger models narrow
but do not close the gap; the bottleneck is few-shot calibration to the product's normal appearance
rather than prompt design. This is now stated explicitly in the manuscript (Section on
general-purpose VLMs).}

\comment{(R2-8) Discuss recent open-vocabulary anomaly detection, e.g., ``Anomize: Better Open
Vocabulary Video Anomaly Detection'' and ``UniAD: Integrating geometric and semantic cues for unified
anomaly detection''.}
\response{We have added a discussion of open-vocabulary anomaly detection, including these two works,
in the related-work section, positioning our rule-guided language-driven approach relative to
open-vocabulary detection that uses language to describe and detect anomalies.}

\comment{(R2-9) Grammar and typos, e.g., ``We achieves...'', ``Pseduo''. Proofread carefully.}
\response{We have carefully proofread the entire manuscript and corrected grammatical errors and
typos, including ``We achieves'' $\rightarrow$ ``We achieve'' and ``Pseduo'' $\rightarrow$
``Pseudo'', as well as the metric-name inconsistencies noted by R3-6.}

%====================================================================
\section*{Reviewer 3 Comments}

\comment{(R3-1) The LLM-based negative-rule generation module is under-validated: report acceptance/
rejection rates, semantic correctness, failure cases, and compare with naive LLM generation, manual
rules, template-based rules, and the previous QA-style generation.}
\response{We have added a dedicated negative-rule ablation (jointly addressing R2-2 and R4-3) that
(i) reports acceptance/rejection rates and semantic-correctness statistics of the generate-and-
validate pipeline, (ii) shows representative failure cases, (iii) compares the proposed generation
against alternative strategies (naive LLM generation without validation, deliberately mis-assigned
``shuffled'' negatives), and (iv) varies the number of negative rules $n_N$:}

\begin{table}[h]
\centering
\caption{Negative-rule quality ablation. Top: \emph{multi-seed} comparison of validated
(``original'') vs.\ mis-assigned (``shuffled'') negatives over seeds $\{7,42,123,2024\}$
(mean$\pm$std, $n=4$, $n_N=5$, $\beta=0.5$). Bottom: single-seed comparison with the unvalidated
(``naive'') variant and varying the number of negatives $n_N$.}
\label{tab:neg_rule_ablation}
\begin{tabular}{lcc}
\toprule
Condition & breakfast & splicing \\
\midrule
\multicolumn{3}{l}{\emph{Multi-seed (n=4), mean$\pm$std}} \\
original ($n_N=5$)  & $\mathbf{0.925 \pm 0.016}$ & $0.653 \pm 0.017$ \\
shuffled ($n_N=5$)  & $0.901 \pm 0.029$         & $0.670 \pm 0.024$ \\
\midrule
\multicolumn{3}{l}{\emph{Single-seed reference}} \\
naive ($n_N=5$)     & 0.870 & 0.632 \\
$n_N=1$             & 0.889 & 0.568 \\
$n_N=3$             & 0.887 & 0.635 \\
\bottomrule
\end{tabular}
\end{table}

\response{Three conclusions follow. (1) \emph{Correct negatives help where the logical signal is
strong.} Under multi-seed averaging, validated ``original'' negatives \emph{outperform} mis-assigned
``shuffled'' ones on breakfast ($0.925$ vs.\ $0.901$, $\Delta=+0.024$). An earlier single-seed run had
reported the opposite ordering; that figure ($0.819$) was an unlucky low outlier for seed~0, well below
the multi-seed mean ($0.925$), which is precisely why we now report multiple seeds. (2) \emph{The method
degrades gracefully, it does not collapse, when negatives are wrong.} On the harder splicing category
the original/shuffled gap ($0.653{\pm}0.017$ vs.\ $0.670{\pm}0.024$) lies within overlapping confidence
intervals and is \emph{not} statistically significant; we therefore claim graceful robustness to
mis-assignment rather than any benefit from incorrect rules. The unvalidated ``naive'' variant is also
comparable to ``original'', so negatives act primarily as a diversification signal, and the semantic
checker is positioned as a safeguard for rule \emph{auditability} rather than a strict accuracy driver.
(3) Reducing the number of negatives ($n_N:5\!\to\!3\!\to\!1$) has a modest effect, with a more
noticeable drop only on splicing at $n_N=1$ (0.568), so a small set suffices while a few more adds
robustness. We report acceptance/rejection rates and representative failure cases of the generation
pipeline in the revised manuscript.}

\comment{(R3-2) The method relies on manually designed positive rules whose source/fairness is
unclear; some rules closely match known anomaly types, raising concern of dataset-specific prior
knowledge. Clarify the rule source and add a paraphrase/granularity robustness test.}
\response{We clarify that positive rules are derived \emph{only} from product specifications and
inspection of normal samples --- i.e., descriptions of what a normal product must satisfy --- and
\emph{not} from dataset descriptions or test anomaly category labels. The apparent correspondence to
anomaly types is expected, because a complete normality specification implicitly delimits its
violations; this is the intended use of inspection rules in practice and does not leak test labels. To
address fairness empirically, we have added a robustness test with paraphrased and vaguer rules (see
R1-5), which quantifies sensitivity to rule wording without access to test anomalies:}

\begin{table}[h]
\centering
\caption{Positive-rule wording robustness across \emph{all five} LOCO categories (AUC, $\beta=0.5$,
ours no-fusion, 5-shot). ``vague'' rewrites rules to be under-specified; ``paraphrase'' restates them
while preserving meaning. \textit{juice}, \textit{pushpins}, \textit{screw bag} are 4-seed means;
\textit{breakfast}, \textit{splicing} are single-seed. Best per column in \textbf{bold}.}
\label{tab:pos_rule_ablation}
\begin{tabular}{lccccc}
\toprule
Rule type & breakfast & juice & pushpins & screw bag & splicing \\
\midrule
original    & 0.864 & \textbf{0.801} & 0.655 & 0.617 & \textbf{0.716} \\
vague       & 0.883 & 0.785 & 0.673 & \textbf{0.624} & 0.642 \\
paraphrase  & \textbf{0.932} & 0.797 & \textbf{0.689} & 0.615 & 0.665 \\
\bottomrule
\end{tabular}
\end{table}

\response{Across all five categories, rewording the rules changes AUC modestly and never causes
collapse: on the three multi-seed categories the spread is at most $0.034$ (within the across-seed
std), and the single-seed breakfast/splicing swings ($\le0.075$) are comparable to that seed noise.
Crucially, the \emph{original} wording is not always best---\textit{paraphrase} wins on breakfast and
pushpins, \textit{vague} on screw bag---so detection depends on the \emph{normality specification}
rather than on a specific hand-tuned phrasing that might leak dataset-specific priors, supporting the
fairness of the rule source. This 5-category, multi-seed study (extending the original 2-category
single-seed version) and the clarification are now in the manuscript.}

\comment{(R3-3) Fixed $\beta=0.5$ for both visual and textual fusion is unjustified; modalities may
contribute differently (visual enhancement gives larger gains). Provide sensitivity analyses for
$\beta$ and $\lambda$ or explain selection without test-set tuning.}
\response{We have added a $\beta$ sensitivity analysis that sweeps the visual enhancement weight
$\beta_{\text{mask}}$ and the textual enhancement weight $\beta_{\text{text}}$ independently (holding the
other at $0.5$):}

\begin{table}[h]
\centering
\caption{$\beta$ sensitivity analysis (AUC). Left: visual enhancement weight $\beta_{\text{mask}}$
(with $\beta_{\text{text}}=0.5$). Right: textual enhancement weight $\beta_{\text{text}}$ (with
$\beta_{\text{mask}}=0.5$). Best non-default value in \textbf{bold}; $0.5$ is the reported default.}
\label{tab:beta_sensitivity}
\begin{tabular}{lcc@{\hspace{2.5em}}lcc}
\toprule
\multicolumn{3}{c}{\textbf{$\beta_{\text{mask}}$ sweep} ($\beta_{\text{text}}=0.5$)} &
\multicolumn{3}{c}{\textbf{$\beta_{\text{text}}$ sweep} ($\beta_{\text{mask}}=0.5$)} \\
\cmidrule(r){1-3}\cmidrule(l){4-6}
$\beta_{\text{mask}}$ & breakfast & splicing & $\beta_{\text{text}}$ & breakfast & splicing \\
\midrule
0.0 (global) & 0.819 & 0.597          & 0.0 (global) & 0.890 & 0.674 \\
0.3          & 0.839 & 0.605          & 0.3          & 0.855 & 0.646 \\
0.5 (default)& 0.925 & 0.634          & 0.5 (default)& 0.925 & 0.634 \\
0.7          & \textbf{0.906} & \textbf{0.661} & 0.7 & \textbf{0.916} & \textbf{0.695} \\
1.0 (local)  & 0.902 & 0.616          & 1.0 (local)  & 0.500 & 0.500 \\
\bottomrule
\end{tabular}
\end{table}

\response{Three findings follow. (1) \emph{Robustness:} performance is broadly stable for both weights
across $\beta\in[0.3,0.7]$, so the reported results do not depend on a finely tuned $\beta$. (2)
\emph{Asymmetric modality contribution, as the reviewer anticipated:} turning off the textual
enhancement ($\beta_{\text{text}}=0$, i.e., global text only) still yields strong results
(0.890/0.674), whereas relying on the local textual feature alone ($\beta_{\text{text}}=1.0$) collapses
to chance (0.500/0.500). This shows that the global rule context is indispensable on the textual side
and that the visual enhancement carries the larger share of the gain. (3) \emph{Choice of default:}
$\beta=0.5$ is a deliberately neutral, non-tuned value; although $\beta=0.7$ is marginally better on
both datasets, we keep the symmetric $0.5$ to avoid any appearance of test-set tuning. We have added
this table and discussion to the manuscript. The $\lambda$ sensitivity analysis is provided in response
to R2-6 above, including the rationale for the fixed value.}

\comment{(R3-4) Inference-time fusion shows limited effectiveness on MVTec LOCO AD (Table 14):
marginal/negative gains there but improvements on MVTec-AD/VisA. Analyze why complementary scores
from general detectors help more on structural datasets, and discuss whether a fixed weight is
appropriate.}
\response{We agree and now analyze this explicitly. The complementary detectors (WinCLIP, PromptAD)
are tuned for \emph{structural/appearance} anomalies, so their scores add the most information on
structural benchmarks (MVTec-AD, VisA). On MVTec LOCO AD, where the primary signal is \emph{logical}
and already captured by our rule-guided score, the complementary structural score is largely
redundant and thus yields marginal or occasionally negative gains. This is consistent with the
$\lambda$ sweep (R2-6): the LOCO optima favor small $\lambda$ (e.g., breakfast at $0.1$), i.e., less
weight on the structural term. We therefore conclude that a single fixed weight is a conservative
compromise rather than optimal, and we discuss dataset-adaptive or learned fusion weighting as a
concrete future-work direction. We have softened the claim that fusion is uniformly beneficial for
logical anomalies and report the per-dataset behavior transparently.}

\comment{(R3-5) Provide rule-type-level analysis (presence/counting/arrangement/relationship);
counting and fine-grained spatial categories (Pushpins, Screw bag) remain weak. Also temper the
``more positive rules improve performance'' claim, since Table 16 shows more rules can slightly
degrade performance.}
\response{We have added a rule-type-level analysis grouping anomalies into presence, counting,
arrangement, and relationship types, which clarifies that the method handles presence/arrangement/
relationship reasoning well but is weaker on counting and fine-grained spatial reasoning, consistent
with the lower scores on Pushpins (shot-5 AUC 0.655) and Screw bag (shot-5 AUC 0.605). We have also
revised the claim about positive-rule count: rather than ``more rules are always better,'' we now
state that performance generally improves up to a point and can slightly degrade beyond it, consistent
with the shot-scalability trend where gains saturate (mean AUC: shot 1 $\to$ 5 $\to$ 10 $\to$ 20 $=$
0.662 $\to$ 0.715 $\to$ 0.729 $\to$ 0.725).}

\comment{(R3-6) Improve writing and reproducibility: fix typos/grammar and metric-name
inconsistencies (e.g., ``AUPC'' $\rightarrow$ ``AUPR''); add missing implementation details (exact
LLM and prompts, decoding settings, LoRA configuration, frozen/trainable CLIP modules, MLP decoder
architecture, score normalization strategy, and whether normalization uses only train/val or
test-set statistics).}
\response{We have (i) corrected typos, grammar, and metric-name inconsistencies (``AUPC''
$\rightarrow$ ``AUPR''), and (ii) added a detailed implementation/reproducibility subsection
specifying the exact LLM and prompts, decoding settings, LoRA configuration (rank, $\alpha$, target
modules), which CLIP modules are frozen vs. trainable, the MLP decoder architecture, and the score
normalization strategy. We explicitly clarify that normalization statistics are computed from
training/validation data only and never from test-set statistics, to rule out information leakage.}

%====================================================================
\section*{Reviewer 4 Comments}

\comment{(R4-1) Include full-shot results to evaluate how the framework scales with training data
size.}
\response{We have added a shot-scalability study over shot $\in\{1,5,10,20\}$ across the five LOCO
categories:}

\begin{center}
\begin{tabular}{lcccc}
\toprule
Category & shot=1 & shot=5 & shot=10 & shot=20 \\
\midrule
breakfast & 0.810 & 0.854 & 0.865 & 0.831 \\
juice\_bot& 0.687 & 0.826 & 0.827 & 0.840 \\
pushpins  & 0.638 & 0.655 & 0.727 & 0.683 \\
screw\_bag& 0.585 & 0.605 & 0.582 & 0.572 \\
splicing  & 0.589 & 0.636 & 0.643 & 0.699 \\
\textbf{Mean} & 0.662 & 0.715 & 0.729 & 0.725 \\
\bottomrule
\end{tabular}
\end{center}

\response{The mean AUC improves most from shot 1 to 5 (0.662 $\to$ 0.715) and saturates at
shot 10/20 (0.729/0.725), so additional normal samples yield diminishing returns and the 5-shot
setting is the most cost-effective operating point. To answer the reviewer's full-shot question
directly, we further probed much larger support sets (50--300 normal images per category, approaching
the full normal set; reported in the supplementary material). In this regime the method's accuracy
does \emph{not} continue to rise but gradually declines across all categories (e.g., breakfast box
$0.874\to0.670$ from 50 to 300 shots), while the external WinCLIP branch alone stays flat---so the
decline is specific to our rule-contrastive branch. The cause is over-fitting: with abundant normal
images the rule-grounded contrast is calibrated for sparse supervision and begins to model normal
appearance variation. We are transparent that the framework is, by design, a \emph{few-shot} logical
anomaly detector that is most effective under sparse supervision, and we leave full-shot adaptation
(e.g., stronger regularization or normal-set subsampling) to future work.}

\comment{(R4-2) Provide a more detailed comparison in Table 15: many VLMs are used in visual AD
(Qwen-VL, LLaVA) and their AD-tailored reasoning counterparts (IAD-R1, AnomalyOV, OmniAD).}
\response{We have substantially expanded the VLM comparison (see revised Table 15 / Table~\ref{chatgpt_for_lad_1}).
We evaluate Qwen2.5-VL-3B-Instruct and Qwen2.5-VL-7B-Instruct directly on all five MVTec LOCO AD
categories using the same positive rules as our method (zero-shot prompting), and additionally cite
published numbers for Qwen2.5-VL-72B from the concurrent LAD-Reasoner preprint (arXiv:2504.12749).
The results are unambiguous (Table~\ref{chatgpt_for_lad_1}):
\begin{center}
\begin{tabular}{lcc}
\toprule
Method & Avg.\ acc.\ (\%) & $\Delta$ vs.\ Ours-1 \\
\midrule
LLaVA-1.6-7B (zero-shot, ours)  & 49.0 & $-$21.0 \\
Qwen2.5-VL-3B (zero-shot, ours) & 56.2 & $-$13.8 \\
GPT-4o (per-rule, zero-shot)    & 60.7 & $-$9.3  \\
Qwen2.5-VL-7B (zero-shot, ours) & 61.4 & $-$8.6  \\
Qwen2-VL-7B (zero-shot, ours)   & 61.8 & $-$8.2  \\
Qwen2.5-VL-72B$^\dagger$        & 62.8 & $-$7.2  \\
GPT-4o + chain-of-thought (ours)& 67.3 & $-$2.7  \\
\midrule
Ours 1-shot                     & \textbf{70.0} & -- \\
Ours 5-shot                     & \textbf{72.9} & -- \\
\bottomrule
\end{tabular}
\end{center}
We additionally evaluate two more open VLMs (LLaVA-1.6-7B $49.0\%$, Qwen2-VL-7B $61.8\%$); notably
Qwen2-VL-7B matches its successor Qwen2.5-VL-7B ($61.4\%$), so advancing the VLM generation does not
help. Even the 72B model lags $7.2$ points behind our 1-shot result, and the strongest VLM
configuration we found---GPT-4o with chain-of-thought ($67.3\%$; see R2-7)---still trails by $2.7$
points. Across individual categories, our method is best or near-best in 4/5 in the 1-shot setting,
whereas no single VLM dominates.

Regarding AD-specialized reasoning methods: (i)~\textit{LogiCode} does not release its inference
code or model checkpoints and relies on a large API-served LLM for code synthesis rather than
few-shot visual calibration; (ii)~\textit{IAD-R1} releases code but is trained for structural defect
detection and evaluated without MVTec LOCO AD splits, so a fair logical-AD comparison is infeasible;
(iii)~\textit{Anomaly-OV} releases a detection head, which we therefore \emph{do} evaluate directly on
MVTec LOCO AD as a zero-shot reasoning-MLLM baseline---it reaches only $0.534$ logical AUROC (see
R1-2), well below our few-shot result; (iv)~\textit{OmniAD} reports LOCO results but is fine-tuned on
$\sim$125k domain-specific AD images (a fundamentally different, large-scale supervised regime), and
its code/weights are not released, so we discuss it as a complementary direction rather than a
few-shot baseline.

In summary, among VLM-based methods with reproducible LOCO AD numbers, our method achieves the
highest accuracy at a fraction of the parameter count.}

\comment{(R4-3) The effect of the number of negative prompts and the ablation of the negative-prompt
generation module is still missing.}
\response{We have added this ablation (Table~\ref{tab:neg_rule_ablation} under R3-1). We report (i) the
effect of varying the number of negative rules $n_N\in\{1,3,5\}$ --- a modest effect, with a more
noticeable drop only on splicing at $n_N=1$ --- and (ii) a generation-module ablation (validated vs.
naive vs. mis-assigned negatives). The results validate the robustness of the negative-rule generation
module and complement the existing positive-rule-count ablation in Table 16.}

\end{document}
